팔레트 자세 추정은 제조·물류 환경에서 화물에 접근하고 포크를 정렬하는 데 필요한 위치와 방향 정보를 제공한다.
그러나 컬러 영상에서 팔레트를 올바르게 검출하더라도 코너 위치와 대응의 오차가 남아 자세 추정의 정확도를 제한할 수 있다.
본 연구는 등록 치수를 사용할 수 있고 초기 팔레트 코너가 이미 산출된 조건에서, 영상 특징과 치수에 따른 국소 키포인트 보정을 평가한다. 검출 실패 복구나 독립 계측 정확도는 평가 범위에 포함하지 않는다.
영상 기반 초기 추정기의 가중치를 고정한 상태에서 보정기를 학습하며, 초기 코너 주변의 다중 해상도 특징과 치수로 이동 후보를 점수화한 뒤 후보 변위의 확률 가중평균으로 좌표를 수정한다.
학습에는 객체의 회전 동등성을 고려한 코너 대응을 사용하고, 보정 전후의 자세 계산에는 동일한 카메라 정보와 치수를 적용한다.
반복 개발에 사용한 직사각형 실사 영상 319장에서 YOLO의 코너 오차 중앙값은 6.72에서 5.78픽셀로, 기하 재구성 참조에 대한 위치·회전 오차 중앙값은 각각 7.90에서 7.07cm, 2.54에서 2.07도로 감소했다.
DOPE와 ResNet-18에도 동일한 구성의 보정기를 각각 학습·적용해 코너 오차 중앙값의 감소를 확인했으며, 별도 정사각형 영상 119장에서도 세 기반의 2차원 보정 효과가 관찰됐다.
다만 큰 오류의 변화와 회전 대응 감독의 추가 이득은 일관되지 않았으며, 결과는 평가한 영상과 참조의 범위에서 국소 보정의 효과를 보여준다.

치수 입력의 추가 효과는 작고 대칭 감독의 변화는 혼합되어 전체 보정의 이득을 두 요소 중 하나의 효과로 설명하지 않는다. PoseFix 방식은 더 낮은 중앙 오차를, N3는 더 낮은 코너 P90을 보였으며, 같은 장치·실행 경계의 시간과 함께 방법 전체의 손익을 보고한다.
